\chapter{Development Environment}
\label{ch:environment}

\section{Tested software stack}
The repository's tested environment is Ubuntu 24.04, ROS 2 Jazzy, Gazebo
Harmonic / Gazebo Sim 8, and system Python 3.12 with NumPy and SciPy. The
validated setup uses the ROS vendor libraries for Gazebo integration and a
matching Gazebo Transport 13 Python binding. These version choices are not
interchangeable details: the Python extension must match the installed
Transport library, and ROS Python packages must be imported by the same Python
interpreter used by the sourced ROS installation.

\begin{description}
 \item[ROS 2 Jazzy] Supplies the node graph, launch system, messages, services,
 actions, parameter handling, and ROS 2 control packages. This project uses
 \code{rclpy}; it has no custom C++ node source.
 \item[Gazebo Harmonic] Simulates rigid-body dynamics, gravity, collisions,
 model poses, the robot, and the detachable-joint plugin. Gazebo Transport is
 also used directly for pose and grasp topics.
 \item[ros\_gz] Connects Gazebo Sim and ROS 2. In the task launch, the explicit
 bridge is the simulation clock. The arm's ros2\_control interface is provided
 by the external \code{gz_ros2_control} system plugin.
 \item[RViz2] Visualizes the URDF and TF tree in the separate display launch.
 The pick-place task does not require or launch RViz.
 \item[Python 3.12, NumPy, SciPy] Implements URDF parsing, transform
 composition, rotation representations, bounded nonlinear least squares,
 constrained far-pickup optimization, sampling, and ROS nodes.
 \item[colcon and ament] Build the three packages in one ROS workspace and
 install package resources and Python modules.
\end{description}

The project does not include source code for a C++ controller or Gazebo system;
those components are installed dependencies. CMake is used to configure and
install the ROS packages, but the project packages primarily install XML,
meshes, launch/configuration files, and Python.

\section{Packages and their responsibilities}
\begin{table}[htbp]
\centering\small
\caption{Three packages built from the \code{gp7_ros2/} workspace.}
\label{tab:packages}
\begin{tabularx}{\textwidth}{l l X X}
\toprule Package & Build type & Main content & Runtime role\\\midrule
\code{gp7_description} & ament CMake & URDF, GP7 STL, gripper DAE, RViz config and display launch & Robot model and visualization resources\\
\code{gp7_gazebo} & ament CMake & SDF world, controller YAML, Gazebo launch & Simulation world and ROS/Gazebo control integration\\
\code{gp7_ros2_pkg} & ament CMake + Python install & \code{gp7_motion}, scripts, launches, tests & Kinematics, planning, control clients, task execution\\
\bottomrule
\end{tabularx}
\end{table}

\section{Workspace and environment sourcing}
A ROS workspace has source packages and generated \code{build/}, \code{install/},
and \code{log/} trees. \code{colcon build} invokes package build systems,
installs resources in the overlay, and records environment hooks. Sourcing
\code{/opt/ros/jazzy/setup.bash} makes the underlay visible; sourcing
\code{install/local_setup.bash} adds this workspace's packages. An overlay
can inadvertently import stale Python or a different ROS build if environment
variables such as \code{PYTHONPATH}, \code{AMENT_PREFIX_PATH}, or
\code{LD_LIBRARY_PATH} point elsewhere.

The validation notes explicitly encountered a Conda Python preceding system
Python and therefore lacking the required NumPy installation. The proven
commands deactivate Conda, clear stale overlay variables, source Jazzy, and
invoke \code{/usr/bin/python3}. The project's helper
\file{gp7_ros2/run_pick_place.sh} also clears some overlay variables and
sources the workspace, but it does not preserve the custom Gazebo Python
binding paths described in the README; direct launch commands are used for
that reason.

\section{Transport binding exception}
The Python task imports \code{gz.transport13} and \code{gz.msgs10}. On the
recorded host, matching apt packages were unavailable, so the Transport 13.5.0
Python extension was built against the ROS vendor library while leaving the
existing C++ library unchanged. The additional Python paths are
\code{\$HOME/.local/gp7-gz-transport13/lib/python} and
\code{/opt/ros/jazzy/opt/gz_msgs_vendor/lib/python}. This is installation
specific, not a universal requirement for every Jazzy/Harmonic machine. The
hash-checked source tag and exact CMake commands are recorded in
\file{docs/setup.md}; chapter 14 summarizes the reproducible setup.

\section{How the tools cooperate}
The package launch reads the installed GP7 URDF and patches package paths and
control configuration into a task description. It starts Gazebo, publishes
the simulated clock to ROS, spawns the model, and activates controllers. The
Python task node requests trajectories from the controller action and
subscribes to feedback. Gazebo's physics system advances the robot and boxes.
The robot-state publisher derives TF transforms from the URDF plus joint
states. RViz can consume that robot description and TF in a separate display
session. Official ROS 2 / Gazebo integration documentation describes the
bridge and controller concepts used here \cite{ros2_jazzy,gazebo_harmonic,ros2_control_jtc}.
